% job it analyst
%%%%%%%%%%%%%%%%%%%%%%%%%%%%%%%%%%%%%%%%%%%%%%%%%%%%%%%%%%%%%%%%%%%%%%%%%%%%%%%%

\pagestyle{highlightmain}
\highlightbar{
    \section{Natixis}
    \textit{09/2013 - 06/2018}\\
    \\Natixis est une banque française de financement et d'investissement proposant des services de gestion d'actifs, d'assurance et de paiement.
    \\Elle appartient au Groupe BPCE, l'un des principaux groupes bancaires européens, qui accompagne des millions de clients via ses réseaux de banque de détail et ses services financiers spécialisés.
    \\\textcolor{accent!90}{\textit{Les équipes de développement logiciel ont depuis été intégrées à BPCE-SI, l'entité du groupe chargée du développement et de la maintenance des applications.}}
    \vspace{0.5em}\\
    \location{Région parisienne, France}
    \homepage{groupebpce.com}{https://www.groupebpce.com/}
    \linkedin{Groupe BPCE}{https://www.linkedin.com/company/bpce}
    \vspace{0.5em}\\
    Référent :\\
    \contact{Benoit JEHAN, Responsable développement}{benoit.jehan@natixis.com}
}
\mainbar{
    \section[\faGears]{Natixis - Ingénieur logiciel}
    \textit{02/2016 - 06/2018}

    Contributions techniques aux projets PGP, BIC et IP de Natixis Payment Solutions : intégration de l'Open Payment Framework (OPF) aux systèmes de paiement et aux échanges interbancaires du Groupe BPCE.

    \subsection{Principales responsabilités}
    \begin{itemize}
        \item Contribution à l'intégration de l'\textcolor{accent}{Open Payment Framework (OPF)} dans les systèmes de paiement du Groupe BPCE pour les échanges internes, les marchés interbancaires et les projets de paiement instantané.
        \\\item Développement et maintenance de traitements batch \textcolor{accent}{Spring Integration} pour les échanges de fichiers de paiement et les intégrations MQ entre les systèmes BPCE et les composants OPF.
        \\\item Contribution aux frameworks d'intégration des flux de paiement PGP, BIC et IP, notamment le déploiement et l'orchestration avec Jenkins et XL Deploy.
        \\\item Participation à la mise en œuvre de services de \textcolor{accent}{paiement instantané}, en lien avec les évolutions réglementaires DSP2 et les traitements bancaires en temps réel.
        \\\item Contribution aux flux d'intégration d'entreprise impliquant \textcolor{accent}{SWIFT}, \textcolor{accent}{SEPA} et les systèmes d'échanges interbancaires.
        \\\item Collaboration avec l'exploitation pour préparer les déploiements, les flux d'intégration et les \textcolor{accent}{mises en production}, en intervenant sur les environnements de développement, de qualification et de tests de performance.
        \\\item Participation aux mises en production et au support opérationnel des intégrations critiques de paiement dans l'écosystème bancaire BPCE.
    \end{itemize}

    \subsection{Environnement technique}
    \tag{Java} \tag{Spring Framework}[integration] \tag{JPA/Hibernate}\\
    \tag{Spring Integration} \tag{Jython} \tag{Maven} \tag{Ant}\\
    \tag{OPF} \tag{SWIFT} \tag{SEPA} \tag{Instant Payment}\\
    \tag{IID} \tag{WebSphere} \tag{Jenkins} \tag{XL Deploy}\\
    \tag{Serena Dimensions} \tag{Oracle} \tag{Eclipse}

    \subsection{Langues utilisées}
    \tag{Français}[quotidien, langue maternelle] \tag{Anglais}[échanges avec le support OPF]
}
\makebody
\clearpage

% job training program
%%%%%%%%%%%%%%%%%%%%%%%%%%%%%%%%%%%%%%%%%%%%%%%%%%%%%%%%%%%%%%%%%%%%%%%%%%%%%%%%

\pagestyle{highlightmain}
\highlightbar{
    \section{Natixis}
    \textit{09/2013 - 06/2018}\\
    \\Natixis est une banque française de financement et d'investissement proposant des services de gestion d'actifs, d'assurance et de paiement.
    \\Elle appartient au Groupe BPCE, l'un des principaux groupes bancaires européens, qui accompagne des millions de clients via ses réseaux de banque de détail et ses services financiers spécialisés.
    \\\textcolor{accent!90}{\textit{Les équipes de développement logiciel ont depuis été intégrées à BPCE-SI, l'entité du groupe chargée du développement et de la maintenance des applications.}}
    \vspace{0.5em}\\
    \location{Région parisienne, France}
    \homepage{groupebpce.com}{https://www.groupebpce.com/}
    \linkedin{Groupe BPCE}{https://www.linkedin.com/company/bpce}
    \vspace{0.5em}\\
    Référent :\\
    \contact{Benoit JEHAN, Responsable développement}{benoit.jehan@natixis.com}

    \section{Clear2Pay (\textit{Éditeur})}
    \textit{10/2015 - 01/2016}\\
    \\Clear2Pay est un éditeur de solutions de paiement basé à Bruxelles, proposant aux établissements financiers des outils de traitement des paiements sécurisés, performants et unifiés à l'échelle internationale.
    \\Son Open Payment Framework et ses logiciels associés permettent aux banques de gérer différents types de paiements et de simplifier les opérations de compensation, de règlement et de paiement à l'échelle de l'entreprise.
    \\\textcolor{accent!90}{\textit{L'entreprise appartient à FIS, un groupe américain de technologies et de services financiers comptant parmi les principaux fournisseurs mondiaux de solutions bancaires et de paiement. FIS accompagne des milliers d'établissements financiers dans le traitement des paiements, les systèmes bancaires, la gestion des risques et de la conformité, ainsi que la banque numérique.}}
    \vspace{0.5em}\\
    \location{Varsovie, Pologne}
    \homepage{fisglobal.com}{https://www.fisglobal.com/}
}
\mainbar{
    \section[\faGears]{Natixis - Formation chez Clear2Pay}
    \textit{10/2015 - 01/2016}

    Programme de formation technique dans le cadre du poste chez Natixis, au sein des équipes de développement FIS/Clear2Pay à Varsovie, en Pologne : Open Payment Framework (OPF), traitements de paiement d'entreprise et pratiques de développement internationales.

    \subsection{Principales responsabilités}
    \begin{itemize}
        \item Intégration directe aux équipes de développement \textcolor{accent}{Clear2Pay} pour contribuer à l'évolution et à la maintenance des composants de paiement OPF.
        \\\item Intervention sur les interfaces OPF, les processus d'orchestration \textcolor{accent}{BPEL} et les intégrations bancaires dans un écosystème de paiement à grande échelle.
        \\\item Contribution au débogage, aux ajustements fonctionnels et à la résolution d'anomalies de localisation logicielle et de flux de traitement des paiements dans les applications d'entreprise.
        \\\item Participation à des processus de développement structurés : workflows Jira, validation technique par l'assurance qualité (QA) et pratiques de livraison fondées sur le \textcolor{accent}{TDD}.
        \\\item Travail quotidien en \textcolor{accent}{anglais} dans une équipe internationale, principalement polonaise, réunissant développeurs, équipes QA, analystes métier et responsables techniques.
        \\\item Acquisition d'une expérience pratique des formats XML bancaires, des processus d'orchestration applicative et des architectures de paiement à grande échelle.
        \\\item Approfondissement des pratiques de génie logiciel au sein d'un éditeur international de solutions de paiement pour les entreprises.
    \end{itemize}

    \subsection{Environnement technique}
    \tag{Java} \tag{BPEL} \tag{Jython} \tag{Gradle}\\
    \tag{OPF} \tag{SWIFT} \tag{SEPA} \tag{Formats XML bancaires}\\
    \tag{IID} \tag{WebSphere} \tag{SVN} \tag{Jira}\\
    \tag{Oracle} \tag{Eclipse}

    \subsection{Langues utilisées}
    \tag{Anglais}[quotidien, professionnel] \tag{Polonais}[quelques mots du quotidien]
}
\makebody
\clearpage

% job web developer & assistant analyst
%%%%%%%%%%%%%%%%%%%%%%%%%%%%%%%%%%%%%%%%%%%%%%%%%%%%%%%%%%%%%%%%%%%%%%%%%%%%%%%%

\pagestyle{highlightmain}
\highlightbar{
    \section{Natixis}
    \textit{09/2013 - 06/2018}\\
    \\Natixis est une banque française de financement et d'investissement proposant des services de gestion d'actifs, d'assurance et de paiement.
    \\Elle appartient au Groupe BPCE, l'un des principaux groupes bancaires européens, qui accompagne des millions de clients via ses réseaux de banque de détail et ses services financiers spécialisés.
    \\\textcolor{accent!90}{\textit{Les équipes de développement logiciel ont depuis été intégrées à BPCE-SI, l'entité du groupe chargée du développement et de la maintenance des applications.}}
    \vspace{0.5em}\\
    \location{Région parisienne, France}
    \homepage{groupebpce.com}{https://www.groupebpce.com/}
    \linkedin{Groupe BPCE}{https://www.linkedin.com/company/bpce}
    \vspace{0.5em}\\
    Référent :\\
    \contact{Benoit JEHAN, Responsable développement}{benoit.jehan@natixis.com}
}
\mainbar{
    \section[\faGears]{Natixis - Développeur logiciel}
    \textit{09/2013 - 09/2015 (alternance)}

    Alternance sur DEFI, plateforme de conservation de titres de Natixis Eurotitres pour la gestion des PEA et des comptes-titres de plusieurs réseaux bancaires, au sein et en dehors du Groupe BPCE.

    \subsection{Principales responsabilités}
    \begin{itemize}
        \item Contribution au développement et à la maintenance d'OIC et de Platine, deux plateformes web en marque blanche destinées respectivement aux clients particuliers et aux agences bancaires.
        \\\item Développement d'interfaces d'administration et d'écrans de support métier avec Java, \textcolor{accent}{JSP}, Hibernate et IBM WebSphere dans un environnement Java d'entreprise.
        \\\item Développement d'interactions front-end et de comportements d'interface avec HTML, CSS, JSP et \textcolor{accent}{jQuery}, dont le glisser-déposer dans les applications destinées aux clients.
        \\\item Mise en place des mappings entre les données Oracle et les objets métier avec Hibernate, en remplacement de traitements SQL manuels.
        \\\item Développement d'outils middleware d'analyse et de visualisation des messages échangés entre les applications Java et les systèmes back-end du \textcolor{accent}{mainframe COBOL DEFI}.
        \\\item Participation à la mise en place de solutions d'analyse et de suivi des usages sur les plateformes destinées aux clients et aux agences.
        \\\item Contribution au support de qualification et à la préparation des déploiements dans des systèmes d'information bancaires à grande échelle.
        \\\item Montée en autonomie rapide sur les développements confiés, dans le cadre de processus de développement structurés.
    \end{itemize}

    \subsection{Environnement technique}
    \tag{Java} \tag{Hibernate} \tag{JSP}[JSTL, taglibs personnalisées]\\
    \tag{HTML} \tag{CSS} \tag{jQuery} \tag{Google Analytics}\\
    \tag{RSA} \tag{WebSphere} \tag{Tomcat}\\
    \tag{Jenkins} \tag{XL Deploy} \tag{Serena Dimensions}\\
    \tag{Oracle} \tag{Eclipse}

    \subsection{Langues utilisées}
    \tag{Français}[quotidien, langue maternelle]
}
\makebody
\clearpage